\documentclass[aspectratio=169,11pt]{beamer}

% Editable Overleaf source. Compile with pdfLaTeX.
% Main talk: 14 slides. Everything after \appendix is backup material.
% The separately required written report is not replaced by these slides.
\usepackage[T1]{fontenc}
\usepackage{graphicx}
\usepackage{booktabs}
\usepackage{tabularx}
\usepackage{amsmath}
\usepackage{amssymb}
\usepackage{hyperref}

% ---------- Edit personal information here ----------
\newcommand{\StudentName}{Zi (Winston) Wang}
\newcommand{\StudentNetID}{ziw5}
\newcommand{\PresentationDate}{October 15, 2026}
\newcommand{\RepositoryURL}{https://github.com/winztonfas/uiuc_commodity_research}
\title[US--China Commodity Prediction]{Predictive Information Across\\US and Chinese Commodity Futures}
\subtitle{Adaptive Cluster Portfolios and Machine Learning}
\author[Zi Wang]{\StudentName\\[3pt]\small NetID: \StudentNetID}
\institute{IE498MLF\\University of Illinois Urbana-Champaign}
\date{\PresentationDate}

% ---------- Style ----------
\definecolor{IllinoisNavy}{HTML}{13294B}
\definecolor{IllinoisOrange}{HTML}{E84A27}
\definecolor{MutedGray}{HTML}{5D6673}
\usetheme{default}
\setbeamercolor{normal text}{fg=IllinoisNavy,bg=white}
\setbeamercolor{title}{fg=IllinoisNavy}
\setbeamercolor{frametitle}{fg=IllinoisNavy}
\setbeamercolor{structure}{fg=IllinoisOrange}
\setbeamercolor{alerted text}{fg=IllinoisOrange}
\setbeamercolor{footline}{fg=MutedGray}
\setbeamerfont{title}{size=\LARGE,series=\bfseries}
\setbeamerfont{frametitle}{size=\Large,series=\bfseries}
\setbeamerfont{subtitle}{size=\large}
\setbeamerfont{footline}{size=\scriptsize}
\setbeamertemplate{navigation symbols}{}
\setbeamertemplate{itemize item}{\small$\bullet$}
\setbeamertemplate{itemize subitem}{\tiny$\bullet$}
\setbeamersize{text margin left=0.65cm,text margin right=0.65cm}
\renewcommand{\arraystretch}{1.15}
\newif\ifbackup
\backupfalse
\setbeamertemplate{footline}{%
  \leavevmode\hbox{%
    \begin{beamercolorbox}[wd=.76\paperwidth,ht=2.6ex,dp=1.3ex,leftskip=.65cm]{footline}%
    IE498MLF\hspace{0.8em}\insertshortauthor
    \end{beamercolorbox}%
    \begin{beamercolorbox}[wd=.24\paperwidth,ht=2.6ex,dp=1.3ex,rightskip=.65cm]{footline}%
    \hfill\ifbackup Appendix \insertframenumber\else\insertframenumber/14\fi
    \end{beamercolorbox}%
  }%
}

% ---------- Figure helper ----------
% Upload figures to figures/ with the exact names used below.
% A placeholder appears if a figure has not been uploaded yet.
\newcommand{\projectfigure}[2]{%
\begin{center}
\IfFileExists{figures/#1}{%
  \includegraphics[width=\linewidth,height=#2,keepaspectratio]{figures/#1}%
}{%
  \IfFileExists{#1}{%
    \includegraphics[width=\linewidth,height=#2,keepaspectratio]{#1}%
  }{%
    \fbox{\parbox[c][#2][c]{0.91\linewidth}{\centering
      \textbf{Upload figure}\\[7pt]\texttt{\detokenize{#1}}\\[7pt]
      \small Place the image in the figures folder.}}%
  }%
}
\end{center}}
\newcommand{\smallnote}[1]{\par\vspace{3pt}{\scriptsize\color{MutedGray}#1\par}}

\begin{document}

% ============================================================
% 1. COVER
% ============================================================
\begin{frame}
\titlepage
\end{frame}

\section{Introduction / Exploratory Data Analysis}

% 2. RESEARCH OBJECTIVE
\begin{frame}{Research objective}
\large
Does information from US commodity futures improve forecasts
of Chinese commodity portfolio returns?
\vspace{0.4cm}
\normalsize
\begin{itemize}
\item Study \textbf{groups of related commodities}, rather than a single futures contract.
\item Construct portfolios whose membership can adapt as historical relationships change.
\item Compare \textbf{domestic-only} models with models that also use US information.
\end{itemize}
\vspace{0.2cm}
\[
\underbrace{\text{Chinese features}}_{\text{historical}}
+\underbrace{\text{US information}}_{\text{available at cutoff}}
\longrightarrow \text{next Chinese portfolio return}
\]
\smallnote{Forecast improvement represents predictive information. It does not establish a causal spillover.}
\note{The motivation is cross-market co-movement. The test is whether US information adds out-of-sample value after controlling for domestic information.}
\end{frame}

% 3. DATA
\begin{frame}{Data and research universe}
\small
\begin{tabularx}{\linewidth}{p{2.5cm}XX}
\toprule
& \textbf{China} & \textbf{United States}\\\midrule
Source & Tushare daily futures data & Databento, CME Globex OHLCV\\
Research window & October 2021--September 2026 & October 2021--September 2026\\
Configured products & 23 & 16\\
Selected-contract rows & 25,843 & 22,251\\
Main fields & OHLC, volume, turnover, open interest & OHLC and volume\\
Date convention & Chinese trading-day labels & UTC calendar-day bars\\\bottomrule
\end{tabularx}
\vspace{0.2cm}
\begin{itemize}
\item Initial sectors: \textbf{metals, energy, agriculture}.
\item Adaptive portfolio anchors: \textbf{gold, copper, crude oil, soybeans}.
\item Chinese wheat remains unavailable because the candidate contracts are illiquid.
\end{itemize}
\smallnote{Counts refer to the configured universe in the research window. First observed selected-contract dates are October 8, 2021 for China and October 3, 2021 for the US. Returns remain in each market's local price currency.}
\note{The Chinese source also stores settlement fields, but this experiment uses close-based returns. Counts are selected product-date rows, not raw expiry contracts.}
\end{frame}

% 4. EDA
\begin{frame}{Exploratory return distributions}
\projectfigure{sector_return_boxplot.png}{4.4cm}
\small
\begin{itemize}
\item Energy displays relatively large return dispersion and extreme observations.
\item Return distributions differ across sectors and markets.
\item Calendar-date comparisons do not establish a tradable lead--lag relationship.
\end{itemize}
\note{Equal-weight sector baskets are transparent descriptive benchmarks. Constituent availability and market calendars differ.}
\end{frame}

\section{Preprocessing, Feature Extraction, Feature Selection}

% 5. PREPROCESSING AND AVAILABILITY
\begin{frame}{Preprocessing and information availability}
\small
\begin{itemize}
\item Construct product returns without treating price jumps between different expiries as returns.
\item Compound US metals weekend returns into Monday. Preserve genuine weekday gaps.
\item Descriptive US copper index: exclude missing members and renormalize available weights.
\item ML targets: \textbf{freeze Chinese contracts and weights at the forecast cutoff}.
\end{itemize}
\vspace{0.1cm}
\textbf{Example: forecasting Monday's full Chinese trading-day return}
\vspace{0.08cm}
\begin{tabularx}{\linewidth}{p{3.5cm}X}
\toprule
Forecast cutoff & Previous Friday, 15:00 China time\\
Target interval & Friday close to Monday close\\
Latest US daily bar & Normally Thursday UTC, complete Friday morning\\
Chinese daily fields & Available according to an assumed publication delay\\\bottomrule
\end{tabularx}
\smallnote{Defaults: Chinese fields become available 24 hours after close. US bars become available 5 minutes after their UTC interval ends. These are assumptions, not measured arrival timestamps.}
\note{Monday's bar can include Friday night. Completed US Friday information therefore overlaps the full Monday target. The China publication-delay assumption introduces additional lag.}
\end{frame}

% 6. CLUSTERS
\begin{frame}{Adaptive cluster portfolios}
\small
\[d_{ij}=\sqrt{\frac{1-\rho_{ij}}{2}}\]
\begin{itemize}
\item Estimate correlations using \textbf{52 completed historical weekly returns}.
\item Use \textbf{average linkage} and distance cut \textbf{0.46}.
\item Retain clusters with at least two products and estimate \textbf{HRP weights monthly}.
\item Singleton anchor: retain previous eligible members and re-estimate their HRP weights.
\end{itemize}
\projectfigure{anchor_cluster_indexes.png}{2.95cm}
\smallnote{Red triangles mark singleton fallback starts. Blue triangles mark effective constituent changes. Local clusters need not match exposures exactly. The ML target additionally enforces cutoff-known weights and contracts.}
\note{Weekly returns are compounded daily returns, not average weekly prices. Historical warmup explains the late start of adaptive portfolios.}
\end{frame}

% 7. FEATURES
\begin{frame}{Feature construction: 39 candidate predictors}
\small
\begin{tabularx}{\linewidth}{Xp{4.2cm}rr}
\toprule
\textbf{Feature family}&\textbf{History / definition}&\textbf{China}&\textbf{US}\\\midrule
Returns, momentum, volatility, shock & Latest/previous returns; 5- and 20-observation windows &8&8\\
Normalized volume & Prior 20-observation mean and maximum &3&3\\
Normalized turnover & Prior 20-observation mean and maximum &3&0\\
Open interest & 1- and 5-observation changes; return interaction &3&0\\
Positive-return breadth & Fraction of active products with positive returns &1&1\\
Portfolio structure & Constituent count and largest weight &2&2\\
Shock--volume interaction & Standardized shock times relative volume &1&1\\\midrule
\textbf{Market-specific total}&&\textbf{21}&\textbf{15}\\\bottomrule
\end{tabularx}
\vspace{0.15cm}
\textbf{Three cross-market features:} 60-observation correlation, 60-observation sensitivity, and standardized 20-observation momentum divergence.
\smallnote{Normalize activity within products before weighted aggregation. Calculate volatility from portfolio returns. Historical activity denominators exclude the current observation. Constant-feature removal occurs inside training folds.}
\note{The exploratory heatmaps do not determine feature selection using the test period. Missing values are handled within training pipelines.}
\end{frame}

\section{Model Fitting and Evaluation}

% 8. MODELS
\begin{frame}{Model fitting and evaluation}
\small
\begin{tabularx}{\linewidth}{p{3cm}X}
\toprule
\textbf{Model}&\textbf{Role}\\\midrule
Ridge regression & Regularized linear relationships between features and returns\\
Random forest & Nonlinear relationships through an average of regression trees\\
Gradient boosting & Sequential trees that correct residual prediction errors\\\bottomrule
\end{tabularx}
\vspace{0.2cm}
\begin{columns}[T]
\begin{column}{0.49\linewidth}
\textbf{Three input sets}
\begin{itemize}
\item China only: 21 features
\item China + US: 36 features
\item China + US + cross: 39 features
\end{itemize}
\end{column}
\begin{column}{0.49\linewidth}
\textbf{Copper evaluation}
\begin{itemize}
\item Training: 686 observations
\item Test: 239 observations
\item Test boundary: October 1, 2025
\end{itemize}
\end{column}
\end{columns}
\[\operatorname{MSE}=\frac{1}{N}\sum_{t=1}^{N}(y_t-\widehat y_t)^2\]
\smallnote{Use identical test dates within each portfolio. Benchmarks predict zero or the training mean. Imputation, constant-feature removal, and scaling use training data only.}
\end{frame}

\section{Hyperparameter Tuning}

% 9. CV AND TUNING
\begin{frame}{Chronological cross-validation and tuning}
\small
\textbf{Copper expanding-window folds}
\begin{tabularx}{\linewidth}{cXXr}
\toprule
Fold&Training ends&Validation period&Train / validation\\\midrule
1&March 28, 2024&April--September 2024&319 / 123\\
2&September 27, 2024&October 2024--March 2025&442 / 118\\
3&March 28, 2025&April--September 2025&560 / 125\\\bottomrule
\end{tabularx}
\vspace{0.17cm}
\textbf{Grid search}
\begin{tabularx}{\linewidth}{p{2.6cm}Xr}
\toprule
Model&Candidates&Combinations\\\midrule
Ridge&$\alpha\in\{0.1,1,10,100,1000\}$&5\\
Random forest&Depth $\in\{3,6\}$; minimum leaf size $\in\{10,30\}$&4\\
Gradient boosting&Trees $\in\{100,200\}$; learning rate $\in\{0.03,0.1\}$; depth $\in\{1,2\}$&8\\\bottomrule
\end{tabularx}
\vspace{0.15cm}
Select the \textbf{lowest mean validation MSE}, then refit on all training observations.
\smallnote{Purge unavailable training labels at validation boundaries. Each portfolio uses 153 CV fits and 9 final refits. Random forest uses 200 trees and feature fraction 0.7. Test outcomes do not select hyperparameters.}
\end{frame}

\section{Ensembling}

% 10. ENSEMBLE
\begin{frame}{Equal-average ensemble}
\[\widehat y_t^{\,\mathrm{ensemble}}=
\frac{\widehat y_t^{\,\mathrm{Ridge}}+\widehat y_t^{\,\mathrm{RF}}+\widehat y_t^{\,\mathrm{GB}}}{3}\]
\begin{itemize}
\item Average predictions, rather than fitted parameters.
\item Construct a separate ensemble for each feature set.
\item Fix weights at one third each; do not optimize on the test.
\end{itemize}
\vspace{0.2cm}
\small
\begin{tabular}{lrr}
\toprule
\textbf{Copper full-feature model}&\textbf{Test RMSE, bp}&\textbf{Direction accuracy}\\\midrule
Ridge&98.26&51.9\%\\
Random forest&97.64&57.3\%\\
Gradient boosting&97.31&56.9\%\\
Ensemble&97.65&56.5\%\\\bottomrule
\end{tabular}
\smallnote{Averaging can reduce model-specific variation, but does not guarantee better forecasts. One basis point equals 0.01 percentage point.}
\end{frame}

\section{Results and Conclusions}

% 11. COPPER RESULTS
\begin{frame}{Copper: small and model-dependent improvement}
\small
\begin{tabularx}{\linewidth}{Xrrrr}
\toprule
&\multicolumn{2}{c}{\textbf{Test RMSE, bp}}&\textbf{MSE reduction}&\textbf{OOS $R^2$}\\
\cmidrule(lr){2-3}
\textbf{Model}&China only&Full&vs.\ China only&vs.\ zero\\\midrule
Ridge&98.05&98.26&$-0.41\%$&$-0.0143$\\
Random forest&97.80&97.64&$+0.32\%$&$-0.0016$\\
Gradient boosting&97.88&97.31&$+1.16\%$&$+0.0051$\\
Ensemble&97.82&97.65&$+0.34\%$&$-0.0018$\\\midrule
Zero forecast&\multicolumn{2}{c}{97.56}&---&0\\
Training-mean forecast&\multicolumn{2}{c}{97.39}&---&$+0.0036$\\\bottomrule
\end{tabularx}
\vspace{0.15cm}
\[R^2_{\mathrm{OOS},0}=1-\frac{\operatorname{MSE}_{\mathrm{model}}}{\operatorname{MSE}_{\mathrm{zero}}}\]
\begin{itemize}
\item Full gradient boosting reduces domestic-model MSE by \textbf{1.16\%}.
\item Its improvement over the zero forecast is only \textbf{0.51\%}.
\item Most configurations do not outperform simple benchmarks.
\end{itemize}
\smallnote{Test: October 9, 2025--September 30, 2026, with 239 observations. Full means China + US + cross-market. Statistical significance has not been established.}
\end{frame}

% 12. FEATURE IMPORTANCE
\begin{frame}{Feature importance: validation permutation analysis}
\small
\textbf{Method:} shuffle a feature in five-observation blocks, then measure validation MSE changes without retraining.
\vspace{0.2cm}
\begin{tabularx}{\linewidth}{Xrr}
\toprule
\textbf{Copper ensemble feature}&\textbf{MSE increase}&\textbf{Positive folds}\\\midrule
Chinese 5-observation momentum&0.187\%&2/3\\
Chinese previous-observation return&0.150\%&3/3\\
Chinese standardized return shock&0.143\%&3/3\\
Chinese latest available return&0.135\%&3/3\\
US previous-observation return&0.124\%&2/3\\\bottomrule
\end{tabularx}
\vspace{0.2cm}
\begin{itemize}
\item Positive importance: shuffling worsens predictions.
\item Negative importance: shuffling improves predictions.
\item Correlated features can substitute for one another.
\end{itemize}
\smallnote{Three chronological validation folds and 10 permutations. Saved CV-selected hyperparameters are reused: this is a non-nested diagnostic. Importance is not causal evidence or a replacement for feature-group comparisons.}
\end{frame}

% 13. OTHER PORTFOLIOS
\begin{frame}{Comparison across four anchored portfolios}
\small
\textbf{Test MSE reduction from adding US + cross-market features}
\vspace{0.2cm}
\begin{tabularx}{\linewidth}{Xrrrr}
\toprule
\textbf{Portfolio}&\textbf{Ridge}&\textbf{RF}&\textbf{Boosting}&\textbf{Ensemble}\\\midrule
Copper&$-0.41\%$&$+0.32\%$&$+1.16\%$&$+0.34\%$\\
Gold&$-0.31\%$&$+0.28\%$&$-0.36\%$&$-0.10\%$\\
Crude oil&$-1.46\%$&$-0.42\%$&$+0.30\%$&$-0.57\%$\\
Soybeans&$-3.62\%$&$-0.33\%$&$-1.37\%$&$-1.53\%$\\\bottomrule
\end{tabularx}
\vspace{0.3cm}
\begin{itemize}
\item Copper provides the clearest initial nonlinear gain.
\item Gold and crude oil show inconsistent incremental value.
\item US features worsen soybean forecasts in this run.
\end{itemize}
\smallnote{Positive values mean improvement over the same model with domestic features. Same protocol for each portfolio. Test samples: copper 239, others 241. Sector comparisons are exploratory.}
\end{frame}

% 14. CONCLUSION -- only one conclusion slide
\begin{frame}{Results and conclusions}
\large
US information provides \alert{limited, model-dependent} forecast improvement in this daily-return baseline.
\vspace{0.25cm}
\normalsize
\begin{itemize}
\item Adaptive clusters describe changing commodity relationships.
\item Cross-market co-movement does not automatically produce useful forecasts.
\item Copper boosting shows a small gain; most models remain close to or below simple benchmarks.
\end{itemize}
\vspace{0.15cm}
\textbf{Limitations and next steps}
\begin{itemize}
\item Publication delays are assumed. Daily bars cannot isolate Chinese daytime returns.
\item Training data are limited and relationships may change across regimes.
\item Future work: intraday alignment, chronological refitting, uncertainty analysis, and regime-conditioned tests.
\end{itemize}
\smallnote{Findings concern predictive information. They do not establish causality or trading profitability.}
\end{frame}

% ============================================================
% APPENDIX -- backup slides, outside the main 12-minute talk
% ============================================================
\appendix
\backuptrue
\section{Appendix}

\begin{frame}{Appendix: code repository and reproducibility}
\textbf{Repository}
\vspace{0.1cm}
\begin{center}\small
\href{\RepositoryURL}{github.com/winztonfas/uiuc\_commodity\_research}
\end{center}
\small
\begin{tabularx}{\linewidth}{p{5cm}X}
\toprule
\textbf{File}&\textbf{Purpose}\\\midrule
\texttt{cluster\_index.py}&Historical clustering and HRP portfolios\\
\texttt{features.py}&Reusable feature calculations\\
\texttt{calculate\_features.py}&Product and portfolio features\\
\texttt{build\_prediction\_data.py}&Availability, frozen targets, and splits\\
\texttt{copper\_ml.py}&Shared single-portfolio experiment\\
\texttt{sector\_ml.py}&Additional portfolios and comparison\\
\texttt{feature\_importance.py}&Validation permutation importance\\
\texttt{plot.py}&Formatted research figures\\\bottomrule
\end{tabularx}
\smallnote{Data and model artifacts are separate from source code. Intermediate and final tables use DuckDB.}
\end{frame}

\begin{frame}{Appendix: configured commodity roots}
\small
\begin{tabularx}{\linewidth}{p{2.1cm}XX}
\toprule
\textbf{Sector}&\textbf{China}&\textbf{US}\\\midrule
Metals&AU, AG, CU, AL, ZN, PB, NI, SN, PT, PD&GC, SI, HG, ALI\\
Energy&SC, FU, LU, BU&CL, HO, RB, NG\\
Agriculture&B, M, Y, C, PM, WH, JR, RI, LH&ZS, ZM, ZL, ZC, ZW, KE, ZR, HE\\\bottomrule
\end{tabularx}
\vspace{0.25cm}
\begin{tabularx}{\linewidth}{p{2.3cm}XX}
\toprule
\textbf{Anchor}&\textbf{China}&\textbf{US}\\\midrule
Gold&AU&GC\\
Copper&CU&HG\\
Crude oil&SC&CL\\
Soybeans&B&ZS\\
Wheat&WH, unavailable portfolio&ZW\\\bottomrule
\end{tabularx}
\smallnote{Anchors identify changing portfolios, not single-contract targets. PT and PD have short histories. Eligibility depends on historical coverage and return variation.}
\end{frame}

\begin{frame}{Appendix: return and volume feature definitions}
\scriptsize
Let $r_t$ be the latest completed portfolio return available at the cutoff.
\begin{tabularx}{\linewidth}{p{4.5cm}X}
\toprule
\textbf{Feature}&\textbf{Definition / window}\\\midrule
Latest return&$r_t$\\
Previous return&$r_{t-1}$\\
Absolute return&$|r_t|$\\
5-observation momentum&$\prod_{j=0}^{4}(1+r_{t-j})-1$\\
20-observation momentum&$\prod_{j=0}^{19}(1+r_{t-j})-1$\\
5-observation volatility&Sample standard deviation of latest 5 returns\\
20-observation volatility&Sample standard deviation of latest 20 returns\\
Standardized shock&$r_t$ / standard deviation of preceding 20 returns, excluding $r_t$\\
Relative volume&$V_t$ / preceding 20-observation mean volume\\
Volume breakout ratio&$V_t$ / preceding 20-observation maximum volume\\
Volume breakout indicator&1 if breakout ratio exceeds 1; otherwise 0\\
Positive-return breadth&Fraction of active constituents with positive returns\\
Shock--volume interaction&Portfolio shock times weighted relative volume\\\bottomrule
\end{tabularx}
\smallnote{Calculate volume ratios per product before weighted aggregation. Windows count observations, not calendar days. Volatility is not annualized.}
\end{frame}

\begin{frame}{Appendix: turnover, open interest, and structure}
\small
\begin{tabularx}{\linewidth}{p{4.7cm}X}
\toprule
\textbf{Feature}&\textbf{Definition}\\\midrule
Relative turnover&Current amount / preceding 20-observation mean amount\\
Turnover breakout ratio&Current amount / preceding 20-observation maximum amount\\
Turnover breakout indicator&1 if current amount exceeds preceding 20-observation maximum\\
1-observation OI change&$\mathrm{OI}_t/\mathrm{OI}_{t-1}-1$\\
5-observation OI change&$\mathrm{OI}_t/\mathrm{OI}_{t-5}-1$\\
Return--OI interaction&Product return times its 1-observation OI change\\
Constituent count&Number of active products\\
Largest weight&Maximum effective portfolio weight\\\bottomrule
\end{tabularx}
\smallnote{Turnover and OI features are available for China. Current activity uses selected-contract data rather than all-expiry totals. OI changes spanning selected-contract rolls are masked. Monetary turnover differs from contract volume.}
\end{frame}

\begin{frame}{Appendix: cross-market features and diagnostics}
\small
\begin{tabularx}{\linewidth}{p{4.5cm}X}
\toprule
\textbf{Feature}&\textbf{Definition}\\\midrule
Rolling correlation&60 forecast observations, requiring at least 30 paired available returns\\
Rolling sensitivity&Rolling Chinese--US covariance / rolling US return variance\\
Momentum divergence&Difference between local 20-observation momenta, each divided by its available 20-observation volatility times $\sqrt{20}$\\\bottomrule
\end{tabularx}
\vspace{0.2cm}
\textbf{Stored diagnostics excluded from the 39-feature allowlist}
\begin{itemize}
\item Publication timestamps and feature ages
\item Fallback, constituent changes, exclusions, and available weight
\item Contract-roll flag, interval bar count, and paired sample count
\end{itemize}
\smallnote{Associations use cutoff-known returns whose endpoints need not coincide. They are not interpreted as contemporaneous correlations or causal coefficients.}
\end{frame}

\begin{frame}{Appendix: selected copper full-model parameters}
\small
\begin{tabularx}{\linewidth}{p{3.1cm}X}
\toprule
\textbf{Model}&\textbf{Selected configuration}\\\midrule
Ridge&$\alpha=1000$\\
Random forest&200 trees; depth 3; minimum leaf size 10; feature fraction 0.7\\
Gradient boosting&100 trees; learning rate 0.03; depth 1; minimum leaf size 20\\\bottomrule
\end{tabularx}
\vspace{0.3cm}
\begin{itemize}
\item Select minimum mean chronological validation MSE.
\item Each model and feature set selects its own configuration.
\item Fit preprocessing inside each training fold.
\item Keep the final pipeline fixed throughout the test period.
\end{itemize}
\end{frame}

\begin{frame}{Appendix: forecast metrics and benchmarks}
\small
\[\mathrm{MSE}=\frac{1}{N}\sum_t(y_t-\widehat y_t)^2,
\qquad \mathrm{RMSE}=\sqrt{\mathrm{MSE}}\]
\[\mathrm{MAE}=\frac{1}{N}\sum_t|y_t-\widehat y_t|\]
\[R^2_{\mathrm{OOS},0}=1-\frac{\sum_t(y_t-\widehat y_t)^2}{\sum_t y_t^2}\]
\[R^2_{\mathrm{OOS},\bar y_{\mathrm{train}}}=1-
\frac{\sum_t(y_t-\widehat y_t)^2}{\sum_t(y_t-\bar y_{\mathrm{train}})^2}\]
\begin{itemize}
\item Zero benchmark: predict 0 every day.
\item Training-mean benchmark: predict the training mean every day.
\item Positive OOS $R^2$: lower MSE than the benchmark.
\item Direction accuracy compares actual and predicted signs.
\end{itemize}
\smallnote{Negative OOS scores indicate benchmark underperformance, not proof that no market signal exists.}
\end{frame}

\begin{frame}{Appendix: boundary purging}
\small
\textbf{First held-out forecast}
\begin{tabularx}{\linewidth}{p{4cm}X}
\toprule
Target trading date&October 9, 2025\\
Forecast cutoff&September 30, 2025, 15:00 China time\\
September 30 target finishes&September 30, 2025, 15:00\\
Assumed label availability&October 1, 2025, 15:00\\\bottomrule
\end{tabularx}
\vspace{0.3cm}
September 30's label is unavailable at the first test cutoff under the assumed publication delay. Exclude that row from training.
\[\text{training label available time}\leq\text{first evaluation cutoff}\]
\smallnote{Apply the same rule within CV. The original rows remain in the database with a purged-boundary flag.}
\end{frame}

\begin{frame}{Appendix: rolling volatility}
\projectfigure{sector_rolling_volatility_20.png}{5.8cm}
\end{frame}

\begin{frame}{Appendix: metals return distribution}
\projectfigure{metal_return_histogram.png}{5.8cm}
\end{frame}

\begin{frame}{Appendix: energy return distribution}
\projectfigure{energy_return_histogram.png}{5.8cm}
\end{frame}

\begin{frame}{Appendix: agriculture return distribution}
\projectfigure{agriculure_return_histogram.png}{5.8cm}
\end{frame}

\begin{frame}{Appendix: clustering at January 2023}
\projectfigure{market_dendrograms_2023-01-01.png}{5.8cm}
\end{frame}

\begin{frame}{Appendix: clustering at January 2026}
\projectfigure{market_dendrograms_2026-01-01.png}{5.8cm}
\end{frame}

\begin{frame}{Appendix: Chinese copper feature correlations}
\projectfigure{china_copper_features.png}{5.8cm}
\end{frame}

\begin{frame}{Appendix: US copper feature correlations}
\projectfigure{us_copper_features.png}{5.8cm}
\end{frame}

\begin{frame}{Appendix: calendar-date cross-market EDA}
\projectfigure{cross_market_copper_calendar_eda.png}{5.8cm}
\smallnote{Calendar-date EDA only. The prediction dataset separately uses availability-based alignment.}
\end{frame}

\begin{frame}{Appendix: copper model comparison}
\projectfigure{copper_ml_comparison.png}{5.8cm}
\end{frame}

\begin{frame}{Appendix: comparison across portfolios}
\projectfigure{sector_ml_summary.png}{5.8cm}
\end{frame}

\begin{frame}{Appendix: copper permutation importance}
\projectfigure{copper_feature_importance.png}{5.8cm}
\smallnote{Validation-only model diagnostic. Repeat variation is not a confidence interval. Group effects are not additive.}
\end{frame}

\begin{frame}{Appendix: joint feature-group importance}
\projectfigure{feature_group_importance.png}{5.8cm}
\smallnote{Validation-only model diagnostic. Repeat variation is not a confidence interval. Group effects are not additive.}
\end{frame}

\begin{frame}{Appendix: data and method references}
\small
\begin{itemize}
\item \href{https://tushare.pro/document/2?doc_id=138}{Tushare: futures daily data and field definitions}
\vspace{0.2cm}
\item \href{https://databento.com/docs/schemas-and-data-formats/ohlcv}{Databento: OHLCV interval conventions}
\vspace{0.2cm}
\item \href{https://www.shfe.cn/eng/reports/CalendarHolidays/TradingHours/}{Shanghai Futures Exchange: trading hours}
\vspace{0.2cm}
\item L\'opez de Prado, M. (2016). \textit{Building Diversified Portfolios that Outperform Out-of-Sample}.\\
\href{https://ssrn.com/abstract=2708678}{HRP methodology}
\vspace{0.2cm}
\item \href{https://scikit-learn.org/stable/modules/permutation_importance.html}{Scikit-learn: permutation feature importance}
\end{itemize}
\end{frame}

\end{document}
